\subsection{Continuous Integration}
\label{subsec:continuous-integration}

Core CI, Profile Matrix, PostgreSQL Integration und Desktop CI prüfen getrennt
Baseline, fünf Profile, PostgreSQL-Kombinationen und native Pakete. Release
Validation läuft nur manuell oder für Versionstags. Seit Governance~v1 besitzt
Core CI bei Pull Request, Push auf \texttt{main} und manueller Ausführung einen
vorgelagerten Ubuntu-Quality-Job mit Python~3.11, Node~20, Rust Stable,
rustfmt, Clippy, Caches und Tauri-Systempaketen. Er ruft \texttt{quality} über
\path{tools/control.py} auf. \texttt{continue-on-error} ist nicht gesetzt;
Tooling, Backend, Frontend und Container deklarieren \texttt{needs: quality}
\parencite{kleiveist2026quality-ci}.

Die Profilmatrix führt das Gate nach Generierung und Installation für alle
fünf Scaffolds vor Tests und Build aus; der generierte PostgreSQL-Pfad tut dies
für drei Backendprofile. Der Release-Workflow setzt Quality vor Tests und
Artefakte. Desktop CI und der direkte PostgreSQL-Job haben dagegen kein eigenes
Gate. Die Reihenfolge
\enquote{Quality vor Tests vor Build} gilt daher nicht pauschal für jeden Workflow.

Die am 9.~August~2026 abgefragten Läufe vom 7.~August gehören zu HEAD
\texttt{a4487ac} und sind nicht vollständig grün. Core CI war für Backend,
Schema, Frontend, Web-Build, Compose und beide Images erfolgreich; Tooling,
Profile und Konfiguration scheiterten. In der Profilmatrix bestanden beide
Webprofile, die drei Desktopprofile scheiterten beim Projekttest. Die
PostgreSQL-Integration war für Master und Web-cloud erfolgreich, nicht für die
beiden Desktopkombinationen. Desktop CI paketierte unter Linux und macOS;
Windows scheiterte bei der
Frontend-Installation
\parencite{kleiveist2026ci-core-run,kleiveist2026ci-profiles-run,kleiveist2026ci-postgres-run,kleiveist2026ci-desktop-run}.

Die am 22.~August~2026 ausgewerteten Läufe für \texttt{461fc751} zeigen ein
erneut gemischtes Bild. Im Core-Lauf war der neue Quality-Job einschließlich
unverändertem Arbeitsbaum erfolgreich; Backend, Frontend und Container
bestanden, während Tooling-, Profil- und Konfigurationstests den
Gesamtworkflow scheitern ließen. In der Profilmatrix bestand
\texttt{web-only}; \texttt{desktop-local} bestand zunächst Quality und
scheiterte später im Projekttest. \texttt{web-cloud},
\texttt{desktop-cloud} und \texttt{full-platform} scheiterten bereits im
Quality-Schritt. Der direkte PostgreSQL-Integrationsjob bestand, die drei
generierten PostgreSQL-Profile scheiterten im Quality-Schritt. Desktop CI
baute unter Linux und macOS, während Windows erneut bei der
Frontend-Installation scheiterte
\parencite{kleiveist2026governance-core-run,kleiveist2026governance-profiles-run,kleiveist2026governance-postgres-run,kleiveist2026governance-desktop-run}.

Öffentliche Metadaten belegen Job und Schrittstatus, die Protokollarchive waren
ohne Berechtigung nicht abrufbar. Die in
Abschnitt~\ref{subsec:systempruefung-installation} beschriebene frische lokale
Reproduktion erklärt die Profilfehler plausibel durch die Ruff-Auflösung; dies
ist eine aus Implementierung, Workflow und Reproduktion abgeleitete Erklärung,
kein direktes Zitat aus den CI-Logs. Der grüne Core-Quality-Job belegt damit
das blockierende Gate im Master-Workflow, nicht die Profilverträglichkeit.

Lokal bestanden am 9.~August~2026 auf \texttt{a4487ac} unter Python~3.13.14 alle
194 Tooling-Tests sowie Schema-, API- und SQLAlchemy-Tests. Das ersetzt weder
die roten externen Jobs noch die mangels Node.js, Rust und Docker nicht
ausgeführten Pfade. Release Validation wurde für diesen Commit nicht
ausgeführt.
\beginnerinput{workspace/statement/600_qualitaet_verifikation/620}
